% !TeX root = ../main.tex
\sloppy
\chapter{Conclusion}
\textit{Mycobacterium tuberculosis} (M.tb) is a clonal and rapidly evolving organism. Current resistant-variant catalogues and the WGS-DST tools built from them (e.g. TB-Profiler) are near static. These tools risk overlooking newly-acquired resistant variants, which lead to False Negative (FN) calls: a susceptible prediction for a resistant sample. This thesis utilised two approaches: first, mine $\sim$45,000 M.tb genomes for variants enriched in FN samples (with True Negative controls), with an emphasis on the PE/PPE gene family, then rank them by independent predictive contribution. Second, assess a single-layer neural network initialised with catalogue weights and then trained as an alternative to static lookup tools. 

The model acted as a proof of concept rather than a polished tool. The single-layer, catalogue-initialised neural network matched or exceeded TB-Profiler sensitivity for 13/15 drugs (14/15 post threshold tuning). Largest gains were observed for data-scarce or newer drugs. Gains were more modest for widely utilised drugs such as the first-line drugs. This sensitivity gain was bought with a specificity floor of 0.80 or 80\%, this is lower compared to the WHO specificity floor of 98\%. The $F_{1}$ score was reduced for lower prevalence drugs. The augmentation layers for the weight schemes brought negligible gains, a consequence of the representational ceiling. Lastly, the trained representations carry predictive value beyond static lookup, and the probabilistic framework can scale which a fixed list cannot do. 

From genome mining, the genuine independent candidate signal was sparse, and concentrated in under-characterisied loci. Post-filtering, only three variants independently moved sensitivity (ablation reduced sensitivity): an $ethA$ variant for CFZ (likely confounder, no CFZ mechanism), $Rv3254$ for CFZ (an unknown hypothetical protein, the strongest genuine candidate) and $pbpB$ for ETO (mechanistically linked, but marginal sensitivity difference). The sparseness of these variants suggests the WHO catalogue already contains independently-predictive canonical variants, with residual signals sitting in auxiliary signals. Ultimately, $Rv3254$ is the top wet-lab validation candidates, with $pbpB$ being a low-confidence lead. 

The central biological contribution came from the PE/PPE gene family. A large set of significant, high-FNP, uncatalogued PE/PPE variants recurred across $\sim$12 drugs and drug classes (notably $PPE54$ \texttt{3733405\_G>A} among first-line drugs). Most parsimonious interpretation suggests most PE/PPE variants were compensatory or resistance-amplifying, not direct drug targets. These variants show no ablation signal as the linear model cannot represent interaction-dependent effects. Therefore, these findings are hypothesis-generating. Overall, this directly answers the PE/PPE research objective, where previously masked regions do harbour FN-associated, uncatalogued variants.

The limitations included data scarcity for newer agents, and no data for certain drugs, which was reflected in their performance in both the model and TB-Profiler. For the model, the single-layer constraint prevented learning of complex resistance mechanisms, as observable in flat augmentation, absent PE/PPE ablation signal, and overall sparse candidate set. 

Future iterations of the model need hidden layers to account for cross-resistance, epistasis and pleiotropy, with a dedicated PE/PPE channel. More resistant data and wet-lab characterisation of $Rv3254$, $pbpB$ and high-prevalence PE/PPE loci would be needed to identify and characterise more candidate variants. This work reframes false-negative reduction as a tractable problem and lays the groundwork for interaction-aware successor model to static catalogues.
